% ============================================================
%  BlueBee 论文主文件
%  编译：latexmk -pdf -interaction=nonstopmode -halt-on-error main.tex
%  产物：build/main.pdf
%  结构说明见 README.md
% ============================================================
\documentclass[journal]{IEEEtran}
% 投会议时改成： \documentclass[conference]{IEEEtran}

\usepackage{amsmath,amssymb}
\usepackage{graphicx}
\usepackage{booktabs}
\usepackage{siunitx}
\usepackage{xcolor}
\usepackage{microtype}
\usepackage{cite}
\usepackage[hidelinks]{hyperref}

% ---- 全局宏：术语、系统名、单位集中在这里改，不要散落正文 ----
\newcommand{\sysname}{BlueBee}
\newcommand{\ble}{BLE}
\newcommand{\stdfour}{IEEE~802.15.4}

\begin{document}

\title{Single-Packet Dual-Protocol Control and Synchronization over BLE Extended Advertising}

% TODO: 填写作者、单位、通讯作者
\author{Author Name
\thanks{Manuscript prepared \today. Funding and affiliation to be completed.}}

\maketitle

\begin{abstract}
% TODO: 150--250 词，按四句写：
%   (1) 问题是什么
%   (2) 本文怎么做（一句话）
%   (3) 关键量化结果（必须与正文表格数字一致）
%   (4) 意味着什么
\emph{Abstract to be written.}
\end{abstract}

\begin{IEEEkeywords}
BLE, IEEE 802.15.4, extended advertising, cross-protocol, synchronization
\end{IEEEkeywords}

\section{Introduction}\label{sec:intro}
% ==== 写作要点 ====
% 1) 背景：BLE 5.0 扩展广播把大负载拆到 secondary；802.15.4/ZigBee 同时工作在 2.4 GHz
% 2) 问题：双协议设备要控制两类接收端时，通常需要发两次
% 3) 本文：同一个 secondary PDU 里同时放 BLE 控制、BlueBee 编码负载和同步字段
% 4) 贡献：3 条，每条后面要能指到具体的节/图/表
% 5) 不要写 "this paper is organized as follows"（IEEE 期刊通常省略）

Bluetooth Low Energy (\ble) extended advertising lets a primary advertisement
reference a secondary payload that is transmitted later, on a channel and at a
time of the sender's choosing~\cite{bluetooth_core51}. The \stdfour{} standard
defines an independent low-rate PHY and MAC that shares the same 2.4\,GHz band
but uses an unrelated frame format~\cite{ieee802154_2020}. Consequently, a node
that must control receivers of both kinds normally transmits twice.

% TODO: 用 2--3 句把"为什么难"讲清楚（不同前导/调制/校验、单包容量、单向广播无确认），
%       再引出本文做法与贡献列表。


\section{Related Work}\label{sec:related}
% ==== 写作要点 ====
% - 按三条线分组：跨协议共存/转换、BLE 扩展广播机制、无线时间同步
% - 每组结尾一句「已有工作没解决什么」，不要逐篇罗列
% - 本节只引已读过的文献，不凭记忆写条目

% TODO: 正文

\section{System Model and Problem Statement}\label{sec:model}
% ==== 写作要点 ====
% - 设备与角色：发射端单板、两类接收端、主机
% - 信道与时频关系：primary/secondary、AUX_ADV_IND 的时频约束
% - 时间基准定义：包参考点、执行时刻、误差如何度量
% - 符号表：本节首次出现的符号全部定义，后续不重复
% - 问题形式化：单包要满足的字节预算与实时约束

% TODO: 正文

\section{Single-Packet Dual-Protocol Design}\label{sec:design}
% ==== 写作要点 ====
% - 字段布局图（一张图顶一千字）：扩展头 / BLE 命令 / BlueBee 区 / 同步字段 / 校验
% - 字节预算表：逐字段开销 + 总计，说明 255 B 上限怎么用
% - 命令语义：run_id、command_id、目标、参数、有效期、去重规则
% - 重发与失效：迟到、过期、单侧漏收时的行为
% - 不要在这里写实测数字，实测放第 VII 节

% TODO: 正文

\section{Dual-Chain Receiver Design}\label{sec:receiver}
% ==== 写作要点 ====
% - 两条独立解码链：BLE 侧与 802.15.4 侧的输入、判决、校验
% - 为什么必须独立解码（不能用一侧结果重建另一侧）
% - 配对与统计：run_id + sequence、完整帧逐字节比对
% - 时间参考估计：结构偏移补偿与残差定义

% TODO: 正文

\section{Implementation}\label{sec:impl}
% ==== 写作要点 ====
% - 硬件：Zynq-7020、AD9363、HackRF、主机
% - 固件：裸机应用的调度、DMA、命令接口
% - 软件：IQ 生成、离线解码、统计脚本
% - 给出可复现所需的最小信息（版本、采样率、增益、带宽）

% TODO: 正文

\section{Evaluation}\label{sec:eval}
% ==== 写作要点 ====
% - 指标定义与口径：分子/分母写清楚，分母为 0 记 N/A
% - 实验设置表：距离、功率、重复轮数、样本量
% - 结果：容量/正确性 → 联合可靠性 → 时延 → 对照（分别发送、无同步补偿）
% - 每张图表必须有正文解读，不能只写「如图所示」
% - 失败样本与异常单独说明，不剔除后只报成功子集

% TODO: 正文

\section{Discussion and Limitations}\label{sec:discussion}
% ==== 写作要点 ====
% - 适用范围：哪些结论只在 SDR 原型上成立
% - 有效性威胁：单 IQ 与双设备、软件事件与执行器动作的区别
% - 尚未验证的部分明确写出来，不要用模糊措辞掩盖

% TODO: 正文

\section{Conclusion}\label{sec:conclusion}
% ==== 写作要点 ====
% - 3--4 句：做了什么、测到什么、边界在哪
% - 不引入正文没有的新数据
% - future work 只写一句，不要展开成第二篇论文

% TODO: 正文


\bibliographystyle{IEEEtran}
\bibliography{refs}

\end{document}
